% Part: first-order-logic
% Chapter: natural-deduction
% Section: proving-things

\documentclass[../../../include/open-logic-section]{subfiles}

\begin{document}

\iftag{FOL}
      {\olfileid[pt]{fol}{ntd}{pro}}
      {\olfileid[pt]{pl}{ntd}{pro}}

\olsection{Exemplos de \usetoken{P}{derivation}}

\begin{ex}
Vamos dar !!a{derivation} da !!{sentence} $(!A \land !B) \lif !A$.

Começamos escrevendo a conclusão desejada na parte inferior da
!!{derivation}.
\begin{prooftree}
\AxiomC{}
\UnaryInfC{$(!A\land !B) \lif !A$}
\end{prooftree}

Em seguida, precisamos descobrir que tipo de inferência poderia resultar em
!!a{sentence} desta forma. O !!{main operator} da
conclusão é $\lif$, então tentaremos chegar à
conclusão usando a regra \Intro{\lif}. O melhor é anotar
as suposições envolvidas e rotular as regras de inferência conforme
avançamos, para que seja fácil ver se todas as suposições foram
!!{discharged}s ao final da prova.
\begin{prooftree}
\AxiomC{$\Discharge{!A \land !B}{1}$}
\DeduceC{$!A$}
\DischargeRule{\Intro{\lif}}{1} 
\UnaryInfC{$(!A\land !B) \lif !A$}
\end{prooftree}

Agora precisamos preencher os passos da suposição $!A \land !B$ até $!A$.
Como só temos um conectivo com o qual lidar, $\land$, devemos
usar a regra $\land$ elim. Isso nos dá a seguinte prova:
\begin{prooftree}
\AxiomC{$\Discharge{!A \land !B}{1}$}
\RightLabel{\Elim{\land}}
\UnaryInfC{$!A$}
\DischargeRule{\Intro{\lif}}{1} 
\UnaryInfC{$(!A\land !B) \lif !A$}
\end{prooftree}
Agora temos uma !!{derivation} correta de $(!A \land
!B) \lif !A$.
\end{ex}

\begin{ex}
Agora vamos dar !!a{derivation} de $(\lnot !A \lor !B)
\lif (!A \lif !B)$.

Começamos escrevendo a conclusão desejada na parte inferior da
!!{derivation}.
\begin{prooftree}
\AxiomC{}
\UnaryInfC{$(\lnot !A \lor !B) \lif (!A \lif !B)$}
\end{prooftree}
Para encontrar uma regra lógica que possa nos dar esta conclusão,
observamos os conectivos lógicos na conclusão: $\lnot$,
$\lor$ e $\lif$. Por enquanto, só nos interessa a primeira
ocorrência de $\lif$, porque é o !!{main operator} da
!!{sentence} na conclusão, enquanto $\lnot$, $\lor$ e a segunda
ocorrência de $\lif$ estão dentro do escopo de outro conectivo, então
cuidaremos deles depois. Começamos, portanto, com a
regra \Intro{\lif}.  Uma aplicação correta deve ser assim:
\begin{prooftree}
\AxiomC{$\Discharge{\lnot !A \lor !B}{1}$}
\DeduceC{$!A \lif !B$}
\DischargeRule{\Intro{\lif}}{1}
\UnaryInfC{$(\lnot !A \lor !B) \lif (!A \lif !B)$}
\end{prooftree}
Isso nos deixa com duas possibilidades para continuar. Podemos
continuar trabalhando de baixo para cima e procurar outra aplicação
da regra \Intro{\lif}, ou podemos trabalhar de cima para baixo e aplicar a
regra \Elim{\lor}. Vamos adotar a segunda opção. Usaremos a suposição
$\lnot !A \lor !B$ como a premissa mais à esquerda de \Elim{\lor}.  Para uma
aplicação válida de \Elim{\lor}, as outras duas premissas devem ser idênticas
à conclusão $!A \lif !B$, mas cada uma pode ser derivada, por sua vez, de
outra suposição, a saber, um dos dois disjuntos de $\lnot !A \lor !B$.
Nossa !!{derivation} ficará, então, assim:
\begin{prooftree}
\AxiomC{$\Discharge{\lnot !A \lor !B}{1}$}
\AxiomC{$\Discharge{\lnot !A}{2}$}
\DeduceC{$!A \lif !B$}
\AxiomC{$\Discharge{!B}{2}$}
\DeduceC{$!A \lif !B$}
\DischargeRule{\Elim{\lor}}{2}
\TrinaryInfC{$!A \lif !B$}
\DischargeRule{\Intro{\lif}}{1} 
\UnaryInfC{$(\lnot !A \lor !B) \lif (!A \lif !B)$}
\end{prooftree}

Em cada um dos dois ramos à direita, queremos derivar $!A
\lif !B$, o que é melhor feito usando \Intro{\lif}.
\begin{prooftree}
\AxiomC{$\Discharge{\lnot !A \lor !B}{1}$}
\AxiomC{$\Discharge{\lnot !A}{2}, \Discharge{!A}{3}$}
\DeduceC{$!B$}
\DischargeRule{\Intro{\lif}}{3}
\UnaryInfC{$!A \lif !B$}
\AxiomC{$\Discharge{!B}{2}, \Discharge{!A}{4}$}
\DeduceC{$!B$}
\DischargeRule{\Intro{\lif}}{4}
\UnaryInfC{$!A \lif !B$}
\DischargeRule{\Elim{\lor}}{2}
\TrinaryInfC{$!A \lif !B$}
\DischargeRule{\Intro{\lif}}{1} 
\UnaryInfC{$(\lnot !A \lor !B) \lif (!A \lif !B)$}
\end{prooftree}

Nas duas partes que faltam na !!{derivation}, precisamos de
!!{derivation}s de $!B$ a partir de $\lnot !A$ e $!A$ no meio, e
a partir de $!A$ e $!B$ à esquerda.  Comecemos pelo primeiro caso. $\lnot
!A$ e $!A$ são as duas premissas de \Elim{\lnot}:
\begin{prooftree}
\AxiomC{$\Discharge{\lnot !A}{2}$}
\AxiomC{$\Discharge{!A}{3}$}
\RightLabel{\Elim{\lnot}}
\BinaryInfC{$\lfalse$}
\DeduceC{$!B$}
\end{prooftree}
Usando \FalseInt, podemos obter $!B$ como conclusão e
completar o ramo.
\begin{prooftree}
\AxiomC{$\Discharge{\lnot !A \lor !B}{1}$}
\AxiomC{$\Discharge{\lnot !A}{2}$}
\AxiomC{$\Discharge{!A}{3}$}
\RightLabel{\Intro{\lfalse}}
\BinaryInfC{$\lfalse$}
\RightLabel{\FalseInt}
\UnaryInfC{$!B$}
\DischargeRule{\Intro{\lif}}{3}
\UnaryInfC{$!A \lif !B$}
\AxiomC{$\Discharge{!B}{2}, \Discharge{!A}{4}$}
\DeduceC{$!B$}
\DischargeRule{\Intro{\lif}}{4}
\UnaryInfC{$!A \lif !B$}
\DischargeRule{\Elim{\lor}}{2}
\TrinaryInfC{$!A \lif !B$}
\DischargeRule{\Intro{\lif}}{1} 
\UnaryInfC{$(\lnot !A \lor !B) \lif (!A \lif !B)$}
\end{prooftree}

Vejamos agora o ramo mais à direita.  Aqui é importante
perceber que a definição de !!{derivation} \emph{permite que hipóteses
  sejam descarregadas} mas \emph{não exige} que o sejam.  Em outras
palavras, se pudermos derivar $!B$ a partir de uma das suposições $!A$ e $!B$
sem usar a outra, tudo bem.  E para derivar $!B$ a partir de~$!B$
é trivial: $!B$ por si só é tal !!a{derivation}, e nenhuma inferência
é necessária.  Portanto, podemos simplesmente eliminar a suposição~$!A$.
\begin{prooftree}
\AxiomC{$\Discharge{\lnot !A \lor !B}{1}$}
\AxiomC{$\Discharge{\lnot !A}{2}$}
\AxiomC{$\Discharge{!A}{3}$}
\RightLabel{\Elim{\lnot}}
\BinaryInfC{$\lfalse$}
\RightLabel{\FalseInt}
\UnaryInfC{$!B$}
\DischargeRule{\Intro{\lif}}{3}
\UnaryInfC{$!A \lif !B$}
\AxiomC{$\Discharge{!B}{2}$}
\RightLabel{\Intro{\lif}}
\UnaryInfC{$!A \lif !B$}
\DischargeRule{\Elim{\lor}}{2}
\TrinaryInfC{$!A \lif !B$}
\DischargeRule{\Intro{\lif}}{1} 
\UnaryInfC{$(\lnot !A \lor !B) \lif (!A \lif !B)$}
\end{prooftree}
Observe que, na !!{derivation} concluída, a inferência \Intro{\lif}
mais à direita na verdade não descarrega nenhuma hipótese.
\end{ex}

\begin{ex}
Até agora não precisamos da regra \FalseCl{}. Ela é especial porque
nos permite descarregar uma hipótese que não é uma sub-!!{formula} da
conclusão da regra.  Está intimamente relacionada à regra \FalseInt{}.
Na verdade, a regra \FalseInt{} é um caso especial da
regra \FalseCl{}---existe uma lógica chamada ``lógica intuicionista'' na
qual apenas \FalseInt{} é permitida.  A regra \FalseCl{} é um último
recurso quando nada mais funciona.  Por exemplo, suponha que queiramos
derivar $!A \lor \lnot !A$. Nossa estratégia usual seria tentar
derivar $!A \lor \lnot !A$ usando $\Intro{\lor}$. Mas isso
exigiria derivar $!A$ ou $\lnot !A$ a partir de nenhuma
suposição, e isso não pode ser feito. \FalseCl{} ao resgate!
\begin{prooftree}
  \AxiomC{$\Discharge{\lnot(!A \lor \lnot !A)}{1}$}
  \DeduceC{$\lfalse$}
  \DischargeRule{\FalseCl}{1}
  \UnaryInfC{$!A \lor \lnot !A$}
\end{prooftree}
Agora procuramos !!a{derivation} de $\lfalse$ a partir de $\lnot(!A \lor
\lnot !A)$. Como $\lfalse$ é a conclusão de $\Elim{\lnot}$, podemos
tentar isso:
\begin{prooftree}
  \AxiomC{$\Discharge{\lnot(!A \lor \lnot !A)}{1}$}
  \DeduceC{$\lnot !A$}
  \AxiomC{$\Discharge{\lnot(!A \lor \lnot !A)}{1}$}
  \DeduceC{$!A$}
  \RightLabel{\Elim{\lnot}}
  \BinaryInfC{$\lfalse$}
  \DischargeRule{\FalseCl}{1}
  \UnaryInfC{$!A \lor \lnot !A$}
\end{prooftree}
Nossa estratégia para encontrar !!a{derivation} de~$\lnot !A$ exige uma
aplicação de~$\Intro{\lnot}$:
\begin{prooftree}
  \AxiomC{$\Discharge{\lnot(!A \lor \lnot !A)}{1}, \Discharge{!A}{2}$}
  \DeduceC{$\lfalse$}
  \DischargeRule{\Intro{\lnot}}{2}
  \UnaryInfC{$\lnot !A$}
  \AxiomC{$\Discharge{\lnot(!A \lor \lnot !A)}{1}$}
  \DeduceC{$!A$}
  \RightLabel{\Elim{\lnot}}
  \BinaryInfC{$\lfalse$}
  \DischargeRule{\FalseCl}{1}
  \UnaryInfC{$!A \lor \lnot !A$}
\end{prooftree}
Aqui, podemos obter $\lfalse$ facilmente aplicando $\Elim{\lnot}$ à
suposição $\lnot(!A \lor \lnot !A)$ e a $!A \lor \lnot !A$, que
segue da nossa nova suposição $!A$ por~$\Intro{\lor}$:
\begin{prooftree}
  \AxiomC{$\Discharge{\lnot(!A \lor \lnot !A)}{1}$}
  \AxiomC{$\Discharge{!A}{2}$}
  \RightLabel{\Intro{\lor}}
  \UnaryInfC{$!A \lor \lnot !A$}
  \RightLabel{\Elim{\lnot}}
  \BinaryInfC{$\lfalse$}
  \DischargeRule{\Intro{\lnot}}{2}
  \UnaryInfC{$\lnot !A$}
  \AxiomC{$\Discharge{\lnot(!A \lor \lnot !A)}{1}$}
  \DeduceC{$!A$}
  \RightLabel{\Elim{\lnot}}
  \BinaryInfC{$\lfalse$}
  \DischargeRule{\FalseCl}{1}
  \UnaryInfC{$!A \lor \lnot !A$}
\end{prooftree}
No lado direito usamos a mesma estratégia, exceto que obtemos $!A$ por~\FalseCl:
\begin{prooftree}
  \AxiomC{$\Discharge{\lnot(!A \lor \lnot !A)}{1}$}
  \AxiomC{$\Discharge{!A}{2}$}
  \RightLabel{\Intro{\lor}}
  \UnaryInfC{$!A \lor \lnot !A$}
  \RightLabel{\Elim{\lnot}}
  \BinaryInfC{$\lfalse$}
  \DischargeRule{\Intro{\lnot}}{2}
  \UnaryInfC{$\lnot !A$}
  \AxiomC{$\Discharge{\lnot(!A \lor \lnot !A)}{1}$}
  \AxiomC{$\Discharge{\lnot !A}{3}$}
  \RightLabel{\Intro{\lor}}
  \UnaryInfC{$!A \lor \lnot !A$}
  \RightLabel{\Elim{\lnot}}
  \BinaryInfC{$\lfalse$}
  \DischargeRule{\FalseCl}{3}
  \UnaryInfC{$!A$}
  \RightLabel{\Elim{\lnot}}
  \BinaryInfC{$\lfalse$}
  \DischargeRule{\FalseCl}{1}
  \UnaryInfC{$!A \lor \lnot !A$}
\end{prooftree}
\end{ex}

\begin{prob}
Dê !!{derivation}s que mostrem o seguinte:
\begin{enumerate}
\item $!A \land (!B \land !C) \Proves (!A \land !B) \land !C$.
\item $!A \lor (!B \lor !C) \Proves (!A \lor !B) \lor !C$.
\item $!A \lif (!B \lif !C) \Proves !B \lif (!A \lif !C)$.
\item $!A \Proves \lnot\lnot !A$.
\end{enumerate}
\end{prob}

\begin{prob}
Dê !!{derivation}s que mostrem o seguinte:
\begin{enumerate}
\item $(!A \lor !B) \lif !C \Proves !A \lif !C$.
\item $(!A \lif !C) \land (!B \lif !C) \Proves (!A \lor !B) \lif !C$.
\item $\Proves \lnot(!A \land \lnot !A)$.
\item $!B \lif !A \Proves \lnot !A \lif \lnot !B$.
\item $\Proves (!A \lif \lnot !A) \lif \lnot !A$.
\item $\Proves \lnot(!A \lif !B) \lif \lnot !B$.
\item $!A \lif !C \Proves \lnot (!A \land \lnot !C)$.
\item $!A \land \lnot !C \Proves \lnot (!A \lif !C)$.
\item $!A \lor !B, \lnot !B \Proves !A$.
\item $\lnot !A \lor \lnot !B \Proves \lnot(!A \land !B)$.
\item $\Proves (\lnot !A \land \lnot !B) \lif\lnot(!A \lor !B)$.
\item $\Proves \lnot(!A \lor !B) \lif (\lnot !A \land \lnot !B)$.
\end{enumerate}
\end{prob}

\begin{prob}
Dê !!{derivation}s que mostrem o seguinte:
\begin{enumerate}
\item $\lnot(!A \lif !B) \Proves !A$.
\item $\lnot(!A \land !B) \Proves \lnot !A \lor \lnot !B$.
\item $!A \lif !B \Proves \lnot !A \lor !B$.
\item $\Proves \lnot \lnot !A \lif !A$.
\item $!A \lif !B, \lnot !A \lif !B \Proves !B$.
\item $(!A \land !B) \lif !C \Proves (!A \lif !C) \lor (!B \lif !C)$.
\item $(!A \lif !B) \lif !A \Proves !A$.
\item $\Proves (!A \lif !B) \lor (!B \lif !C)$.
\end{enumerate}
(Todos exigem a regra $\FalseCl$~.)
\end{prob}

\end{document}
